\section{Introduction}

\subsection{Physical requirements for autonomous reaction mechanism research}

Scientific validation of an organic reaction mechanism requires checking the molecular geometries, reaction pathways, and quantum-chemical results assigned to each species and elementary step. Normal termination of a geometry optimization (exit 0) does not establish that the result is a stable intermediate, and one imaginary frequency does not by itself identify the saddle point of the intended reaction. An intrinsic reaction coordinate (IRC) that reaches an unintended dissociation product or stereoisomer also changes the mechanistic assignment of the associated energy calculation. These checks are particularly important when a study includes several elementary steps or competing stereochemical pathways.\citep{ref09,ref10}

Computational chemistry agents automate several parts of this work. El Agente Q \citep{ref01}, Aitomia \citep{ref02}, and Gráfico \citep{ref06} orchestrate quantum-chemical calculations and typed workflows. Estructural \citep{ref04} and ARCHE \citep{ref05} provide tools for molecular fragment editing, atom mapping, and transition-state geometry preparation; Quntur \citep{ref03} addresses reaction kinetics and transition-state research. Autonomous mechanism exploration also requires reliable geometric operations and a completion decision based on scientific validation.

Aligning stereochemical orientations, assembling reactant fragments, and preparing pathway endpoints require three-dimensional geometric reasoning. Language models can make errors when they infer spatial relationships or directly manipulate Cartesian coordinates. Quntur \citep{ref03}, for example, reports the need for human assistance with endpoint alignment in complex pathway and transition-state preparation. This motivates interfaces through which a model can specify a chemical objective and inspect the resulting geometric change.

Completion requires checking whether the computed species and steps satisfy the physical acceptance conditions. Intermediates require confirmation of zero imaginary frequencies, while the frequency and bidirectional IRC results for a saddle point must correspond to the same reaction. An agent can otherwise declare a mechanism complete while these checks remain unresolved. A system for autonomous exploration therefore needs to retain the chemist's research objective, the acceptance requirements, and the computational artifacts used to evaluate them.

\subsection{Cloudra: a physical contract-driven system for autonomous mechanism exploration}

We introduce Cloudra, a multi-agent system for exploring reaction mechanisms involving small organic molecules. Language models propose mechanistic hypotheses, choose exploration pathways, and assess computational setbacks. Deterministic software maintains atom-topology correspondences and performs local geometric operations. Physical acceptance checks evaluate whether the returned structures, pathways, and mechanistic conclusions meet the task's scientific requirements.

Figure~\ref{fig:01} summarizes a structured literature sample of 85 chemical AI agent systems: 39 concern materials or surfaces, 23 concern finite molecules or complexes, 12 concern soft matter or biomolecules, and 11 cover multiple categories of objects. Figures~\ref{fig:02} and \ref{fig:03} show the documented tasks in the finite-molecule/complex category and compare Cloudra's capabilities with those of the existing systems. The accompanying literature materials contain the complete matrix and the sources for each entry.

\begin{figure}[tbp]
\centering
\includegraphics[width=\linewidth]{figures/fig01.pdf}
\caption{Physical targets in a structured literature sample of 85 chemical AI agent systems: 39 materials/surface, 23 finite-molecule/complex, 12 soft-matter/biomolecular, and 11 cross-class systems.}
\label{fig:01}
\end{figure}

Cloudra uses the principle of intelligent restraint to assign geometric calculations and coordinate changes to deterministic operators. A maximum common-subgraph (MCS) algorithm retrieves molecular fragments. Atom mapping and reaction-center identification connect molecular topology to the coordinates used in subsequent operations. Measurement operators return bond lengths, bond angles, dihedral angles, and stereochemical features; manipulation operators perform fragment rotation, translation, and transition-state assembly. The model specifies the intended structural transformation, examines the measured change, and selects the next configuration. This interface reduces the need for manual endpoint alignment.

\begin{figure}[tbp]
\centering
\includegraphics[width=\linewidth]{figures/fig02.pdf}
\caption{Documented tasks in the finite-molecule/complex class (panel 1 of 2; 23 chemical AI agent systems in total). T: structure construction/editing; E: static energy/properties; O: geometry optimization; F: vibrations/thermochemistry; R: explicit pathways/TS; D: molecular/quantum dynamics. Filled cells denote positive evidence; blanks are unencoded, not absent capabilities.}
\label{fig:02}
\end{figure}

Cloudra organizes scientific validation through a directed acyclic graph centered on chemical entities, the Delivery DAG. Ground-state chemical species and elementary reaction steps provide the references for the associated computational artifacts. Each step carries requirements for a first-order saddle point with exactly one valid imaginary frequency and for bidirectional IRC evidence. The graph binds each claim to the original quantum-chemical logs used in its evaluation. The Supervisor reviews the results against the task's scientific contracts alongside the programmatic acceptance checks. A step with unresolved frequency or pathway requirements remains unconfirmed when the model requests completion.

\begin{figure}[tbp]
\centering
\includegraphics[width=\linewidth]{figures/fig03.pdf}
\caption{Chemical AI agent systems in the finite-molecule/complex class, continued (panel 2 of 2), using the codes in Figure~\ref{fig:02}. Cloudra (this work) is benchmarked at the bottom. Complete matrices for all 85 chemical AI agent systems and source-level qualifications accompany the draft.}
\label{fig:03}
\end{figure}

\subsection{Contributions}

The Delivery DAG assigns scientific validation requirements to chemical species and elementary steps and records which requirements remain unresolved. The checks use the underlying quantum-chemical artifacts as their basis. This allows the system to evaluate completion requests against the physical acceptance conditions throughout an extended research trajectory.

Cloudra also implements molecular geometric operations as deterministic procedures, including MCS-based topological mapping, steric graph traversal, and transition-state fragment assembly. Language models use these procedures to investigate chemical hypotheses, compare competing pathways, and revise unsuccessful calculations during autonomous transition-state searches.

We evaluate these choices with CloudraBench, which measures graph coverage, physical self-consistency, and quantitative barrier accuracy from native computational outputs. Its 50 tasks cover elementary and multistep reactions, stereochemistry, and the applicability limits of the allowed methods. The experiments comprise 48 independent runs with three repeats on four tasks and a completed study of 200 runs under four workflow configurations across all 50 tasks. Standard uses the same model backbone as the primary comparison; Record and Brief-errors examine evidence organization and diagnostic feedback.
